\documentclass[11pt, letterpaper]{article}
\usepackage[utf8]{inputenc}
\usepackage[T1]{fontenc}
\usepackage{amsmath}
\usepackage{amsfonts}
\usepackage{amssymb}
\usepackage{geometry}
\usepackage{verbatim}

\geometry{letterpaper, margin=1in}

\title{\textbf{Problem A: Touko's Perfection}}
\author{Time Limit: 2.0s | Memory Limit: 256MB}
\date{}

\begin{document}

\maketitle

\section*{Description}

"I have to be perfect. Everyone expects it."

Touko Nanami is obsessing over the Student Council's workflow. She has quantified the "Competence Level" of the $N$ members of the extended council (including volunteers) as a sequence $A_1, A_2, \dots, A_N$.

For the workflow to be seamless, the competence levels of the members must be **non-decreasing**. That is, the competence of the first member must be less than or equal to the second, and so on ($B_1 \le B_2 \le \dots \le B_N$).

Touko can "train" any member to change their competence level from $A_i$ to any integer $B_i$. The effort required to change a member's competence is $|A_i - B_i|$.

Touko wants to minimize the **total effort** $\sum_{i=1}^{N} |A_i - B_i|$ required to make the sequence non-decreasing.

\section*{Input}

The first line contains a single integer $N$ ($1 \le N \le 2 \cdot 10^5$).
The second line contains $N$ integers $A_1, A_2, \dots, A_N$ ($1 \le A_i \le 10^9$).

\section*{Output}

Output the minimum total effort required.

\section*{Examples}

\subsection*{Example 1}
\textbf{Input:}
\begin{verbatim}
5
3 2 5 1 7
\end{verbatim}

\textbf{Output:}
\begin{verbatim}
5
\end{verbatim}

\textbf{Explanation:}
We can modify the sequence to $2, 2, 5, 5, 7$.
Cost: $|3-2| + |2-2| + |5-5| + |1-5| + |7-7| = 1 + 0 + 0 + 4 + 0 = 5$.
Alternatively, $2, 2, 2, 5, 7$ costs $|3-2| + |2-2| + |5-2| + |1-5| + |7-7| = 1 + 0 + 3 + 4 + 0 = 8$, which is worse.

\subsection*{Example 2}
\textbf{Input:}
\begin{verbatim}
4
1 1 1 1
\end{verbatim}

\textbf{Output:}
\begin{verbatim}
0
\end{verbatim}

\end{document}
